\apendice{Formalização complementar das soluções do Sudoku}

Este apêndice reúne a parte mais abstrata da formalização do Z3. Ela detalha a ideia apresentada na metodologia: o programa enumera soluções completas para descobrir quais valores continuam possíveis em cada célula.

Seja $\Phi_G$ o conjunto de restrições de uma grade $G$. Uma solução completa é uma atribuição $M$ de valores às células que satisfaz todas essas restrições. O conjunto de soluções é:

\[
\operatorname{Sol}(G)=\{M : M\models\Phi_G\}.
\]

Nessa expressão, $M\models\Phi_G$ significa que o preenchimento $M$ respeita as regras da grade, as pistas e as restrições de linha, coluna e bloco.

Para uma célula localizada na posição $[i,j]$, os valores que aparecem entre todas as soluções completas são definidos por:

\[
\operatorname{Poss}_G(i,j)=\{M(x_{ij}) : M\in\operatorname{Sol}(G)\}.
\]

Se $\left|\operatorname{Poss}_G(i,j)\right|=1$, apenas um valor aparece naquela posição e a célula é globalmente forçada pela grade. Se $\left|\operatorname{Poss}_G(i,j)\right|>1$, mais de um valor pode ocupar a posição em soluções diferentes. O código usa essa distinção ao avaliar se a jogada proposta pela LLM é apenas compatível com alguma solução ou se é única naquela célula.
